\section{MARKED-CUTSET-BLOCK1: a nonvacuous PRE-A member}
\label{sec:construction-strengthen6-wrapper}

An exact partial member now realizes the relaxed WRAPPER ledger.  In a
canonical binary CNF encoding, tautological source markers specify a cutset
$S$ and ordered blocks with disjoint non-$S$ supports.  Before any child SAT
call, the member verifies the raw marker grammar, polynomially bounded
identifiers, all supports, all
$2^{|S|}$ restrictions, and every complete child encoding.  It accepts only
when
\[
 |S|\le\tfrac1{10}\log_2N,\quad b\le N^{1/2},\quad
 |F_{a,j}|\le N^{1/2}\log N.
\]
It then streams the children and evaluates OR over cut assignments and AND
over residual blocks.  Exactness follows from disjoint residual supports.

Counting identical executions separately, the member makes at most
$2^{|S|}b\le N^{3/5}$ calls.  Each source record enters at most one block per
assignment, while complete headers cost $O(b\log N)$, so aggregate child
volume is $N^{11/10+o(1)}$.  The full source-only dry run and emission replay
cost $N^{11/10}\operatorname{polylog}N\le N^{6/5+o(1)}$.  Nonchild retained
state is polylogarithmic and only one $N^{1/2+o(1)}$ child buffer is live.
Malformed and late-failing annotations return UNKNOWN with zero children but
their full discovery cost is charged.

The logical construction is the one-level supplied-certificate case of
recursive/component backdoors; it is not a new backdoor theorem.  The advance
is the causal source implementation and exact encoded-bit ledger.  It does
not find annotations on general CNFs, and arbitrary emitted children need not
be recursively annotated.  Therefore it neither covers the guarded general
core nor eliminates the inherited $K$.  Conditional $K$ elimination remains
valid only under the previously stated all-input recursive closure.  P versus
NP remains OPEN.
